\clearpage\pagestyle{otMain}\setlength{\otBlocEspacement}{2.5mm}\fontsize{11}{15}\selectfont\corrigesChapitre{17}{La symétrie axiale}
\begin{blocLivre}{cor-17-1}\begin{corrige}{1}{Axe ou pas ?}\demarche{Imaginer le pliage le long de la droite.}
Le pliage superpose les deux moitiés du cœur et du sapin : $\res{\text{oui}}$ pour les deux. Pour le drapeau, la partie gauche déborde après pliage : $\res{\text{non}}$.
\label{fig-cor-17-1}\figureGeo[corrige]{sym-1}
\end{corrige}\end{blocLivre}
\begin{blocLivre}{cor-17-2}\begin{corrige}{2}{Compter les axes}\demarche{Tester médianes et diagonales séparément.}
Le carré en a $\res4$ : deux médianes et deux diagonales. Le rectangle en a $\res2$ : ses diagonales n’en sont pas. Le triangle équilatéral en a $\res3$.
\label{fig-cor-17-2}\figureGeo[corrige]{sym-2}
\end{corrige}\end{blocLivre}
\begin{blocLivre}{cor-17-3}\begin{corrige}{3}{La diagonale piégée}\demarche{Distinguer « même forme » et « superposés par pliage ».}
Les deux triangles sont bien superposables, mais le pliage ne les superpose pas : des coins dépassent. La diagonale $\res{\text{n’est pas un axe}}$ de symétrie : l’axe exige la superposition par pliage.
\label{fig-cor-17-3}\figureGeo[corrige]{sym-3}
\end{corrige}\end{blocLivre}
\begin{blocLivre}{cor-17-4}\begin{corrige}{4}{Le symétrique d’un point}\demarche{Compter les carreaux jusqu’à l’axe, puis autant au-delà.}
A est à 2 carreaux à gauche de l’axe : son symétrique est à 2 carreaux à droite, sur la même ligne. De même pour B (3 carreaux) et C (1 carreau). Les réponses figurent en vert.
\label{fig-cor-17-4}\figureGeo[corrige]{sym-4}
\end{corrige}\end{blocLivre}
\begin{blocLivre}{cor-17-5}\begin{corrige}{5}{Compléter la maison}\demarche{Traiter sommet par sommet avant de relier.}
Chaque sommet reçoit son symétrique de l’autre côté, à même distance de l’axe ; on relie ensuite dans le même ordre. La moitié ajoutée figure en vert.
\label{fig-cor-17-5}\figureGeo[corrige]{sym-5}
\end{corrige}\end{blocLivre}
\begin{blocLivre}{cor-17-6}\begin{corrige}{6}{Axe horizontal}\demarche{Compter au-dessus puis au-dessous.}
A est à 1 carreau au-dessus de l’axe : son image est à 1 carreau au-dessous. B, à 3 carreaux, descend à 3 carreaux. Le segment image, en vert, mesure autant que [AB].
\label{fig-cor-17-6}\figureGeo[corrige]{sym-6}
\end{corrige}\end{blocLivre}
\begin{blocLivre}{cor-17-7}\begin{corrige}{7}{Le papillon}\demarche{Les points de l’axe restent en place.}
Les sommets de l’aile dessinée se reflètent sous l’axe, à égale distance. Les points posés sur l’axe ne bougent pas. L’aile ajoutée figure en vert.
\label{fig-cor-17-7}\figureGeo[corrige]{sym-7}
\end{corrige}\end{blocLivre}
\begin{blocLivre}{cor-17-8}\begin{corrige}{8}{Les lettres}\demarche{Tester les deux directions de pliage.}
A : axe $\res{\text{vertical}}$. B : axe $\res{\text{horizontal}}$. H : $\res{\text{les deux}}$. F : $\res{\text{aucun}}$.
\label{fig-cor-17-8}\figureGeo[corrige]{sym-8}
\end{corrige}\end{blocLivre}
\begin{blocLivre}{cor-17-9}\begin{corrige}{9}{Chercher l’erreur}\demarche{Comparer les distances à l’axe, point par point.}
Le sommet $\res{B'}$ est à 2 carreaux de l’axe alors que B est à 3 carreaux : les distances à l’axe doivent être égales. Le point corrigé figure en vert.
\label{fig-cor-17-9}\figureGeo[corrige]{sym-9}
\end{corrige}\end{blocLivre}
\begin{blocLivre}{cor-17-10}\begin{corrige}{10}{Sur l’axe}\demarche{Un point de l’axe ne bouge pas.}
A est sur l’axe : il est $\res{\text{son propre symétrique}}$. Seul B se reflète ; le segment image part donc du même point A. Le tracé figure en vert.
\label{fig-cor-17-10}\figureGeo[corrige]{sym-10}
\end{corrige}\end{blocLivre}
\begin{blocLivre}{cor-17-11}\begin{corrige}{11}{Le carrelage}\demarche{Associer les cases deux par deux à travers l’axe.}
Chaque case colorée a sa jumelle de l’autre côté, à même distance de l’axe et sur la même ligne : 6 cases sont à colorier. Elles figurent en vert.
\label{fig-cor-17-11}\figureGeo[corrige]{sym-11}
\end{corrige}\end{blocLivre}
\begin{blocLivre}{cor-17-12}\begin{corrige}{12}{Deux axes à la fois}\demarche{Appliquer les deux pliages l’un après l’autre.}
On reflète d’abord à travers l’axe vertical, puis tout le résultat à travers l’axe horizontal : la figure finale compte 4 quarts identiques. Les trois quarts ajoutés figurent en vert.
\label{fig-cor-17-12}\figureGeo[corrige]{sym-12}
\end{corrige}\end{blocLivre}
\begin{blocLivre}{cor-17-13}\begin{corrige}{13}{Les longueurs conservées}\demarche{Citer la propriété de conservation.}
La symétrie conserve les longueurs : le segment image mesure $\res4$ carreaux, comme [CD]. La construction en vert le confirme.
\label{fig-cor-17-13}\figureGeo[corrige]{sym-13}
\end{corrige}\end{blocLivre}
\begin{blocLivre}{cor-17-14}\begin{corrige}{14}{Le bateau du port}\demarche{Commencer par les points les plus proches de l’axe.}
Coque et voile se reflètent sommet par sommet ; le mât, posé sur l’axe, reste en place. Le demi-dessin ajouté figure en vert.
\label{fig-cor-17-14}\figureGeo[corrige]{sym-14}
\end{corrige}\end{blocLivre}
\begin{blocLivre}{cor-17-15}\begin{corrige}{15}{Vrai ou faux ?}\demarche{Chercher les cas particuliers avant de valider.}
La première est $\res{\text{presque vraie}}$ : elle oublie les points situés sur l’axe, qui ne bougent pas. La seconde est $\res{\text{fausse}}$ : le carré a 4 axes et le cercle en a une infinité.
\label{fig-cor-17-15}\figureGeo[corrige]{sym-15}
\end{corrige}\end{blocLivre}
